\section{Cronograma}

El cronograma se plantea para el semestre II/2026, del 10 de agosto al 26 de diciembre de 2026. Las actividades se distribuyen en cinco meses, de acuerdo con el diseño metodológico.

\small
\begin{tabularx}{\textwidth}{|p{0.08\textwidth}|X|c|c|c|c|c|}
\hline
\textbf{OE} & \textbf{Actividades} & \textbf{Ago.} & \textbf{Sep.} & \textbf{Oct.} & \textbf{Nov.} & \textbf{Dic.} \\
\hline
1 & Definir escenarios de amenazas digitales y flujo de niveles. & X & X & & & \\
\hline
2 & Diseñar mecánicas de toma de decisiones y consecuencias. & & X & X & & \\
\hline
3 & Organizar contenidos de ciberseguridad por dificultad progresiva. & & X & & & \\
\hline
4 & Construir casos prácticos de \textit{phishing}, malware e ingeniería social. & & & X & X & \\
\hline
5 & Implementar reportes de retroalimentación por escenario. & & & & X & \\
\hline
6 & Realizar pruebas de usabilidad y ajustes del prototipo. & & & & & X \\
\hline
7 & Redactar conclusiones, documentación final y revisión del proyecto. & & & & & X \\
\hline
\end{tabularx}
\normalsize
